\section{Original complex-submanifold charts and corestriction}

\begin{lemma}[A closed complex submanifold is a complex manifold, and its inclusion an immersion]
\label{ana:submanifold-complex-structure}
\uses{ana:geometric-setting,ana:blowups}
\lean{Hironaka.SubmanifoldChartedSpace.chartedSpace,Hironaka.SubmanifoldChartedSpace.isComplexManifold_chartedSpace,Hironaka.SubmanifoldChartedSpace.contMDiff_val}\upstreamlean{TranscendentalNakamaye.Pending.StepFourOnStrictTransform.fderiv_coordinateMap_val_injective}
\mixedleanok
Let \(W\) be a subset of a complex manifold \(X\) of dimension \(n\) which is regular of
dimension \(k\le n\) at each of its points. With the subspace topology and the atlas of
flattening charts (at each point, a holomorphic chart of \(X\) followed by a biholomorphism
straightening \(W\) to \(\mathbf C^k\times0\), restricted to \(W\)), \(W\) is a complex manifold
of dimension \(k\), and the inclusion \(W\hookrightarrow X\) is holomorphic, with injective
differential in every pair of charts.
\end{lemma}
\begin{proof}
\mixedleanok
The transition between two flattening charts is the composite of holomorphic maps restricted
to \(\mathbf C^k\times0\), hence holomorphic; Hausdorffness and second countability pass to
the subspace. In a flattening chart \(F\) and the ambient chart it is built from, the inclusion
reads \(\zeta\mapsto\Phi_F^{-1}(\zeta,0)\), which is holomorphic. In a flattening chart \(F\) and an
ambient chart \(d\), the inclusion read in coordinates has the holomorphic left inverse \(w\mapsto\operatorname{pr}_k(F(d^{-1}(w)))\)
near the image point, so its differential is injective.
\end{proof}

\begin{lemma}[Holomorphic maps into a submanifold]
\label{ana:submanifold-corestriction}
\uses{ana:submanifold-complex-structure}
\source{transcendental-nakamaye-copy4-20260911:prop:principalization}
\lean{Hironaka.SubmanifoldCodRestrict.contMDiff_codRestrict,Hironaka.SubmanifoldCodRestrict.contMDiff_restrictPreimage}
\leanok
Let \(W\subseteq X\) be as in Lemma~\ref{ana:submanifold-complex-structure}, with its complex
structure. If \(f:M\to X\) is a holomorphic map of complex manifolds with \(f(M)\subseteq W\), then
\(f\), regarded as a map \(M\to W\), is holomorphic. In particular, if \(\pi:Y\to X\) is holomorphic
and \(W'\subseteq Y\) is regular of some dimension \(k'\) at each of its points, with
\(\pi(W')\subseteq W\), then the restriction \(\pi|_{W'}:W'\to W\) is holomorphic.
\end{lemma}
\begin{proof}
\leanok
Let \(y\in M\) and let \(F\) be the flattening chart of \(W\) at \(f(y)\): a chart \(e_F\) of \(X\)
followed by a biholomorphism \(\Phi_F\) straightening \(W\) to \(\mathbf C^k\times0\). The map
\(f\) is continuous into \(W\) for the subspace topology, and in the chart \(F\) it reads
\(z\mapsto\operatorname{pr}_k(\Phi_F(e_F(f(z))))\) near \(y\), a composite of holomorphic maps. For the
second assertion apply the first to \(\pi\circ\iota_{W'}\), where the inclusion
\(\iota_{W'}:W'\to Y\) is holomorphic by Lemma~\ref{ana:submanifold-complex-structure}.
\end{proof}
